\begin{recipe}
[% 
    preparationtime = {\unit[10]{min} + \unit[12]{h} Ziehzeit},
    portion = {\portion{4}},
    source = {\protect\recipesource{https://www.coffeecircle.com/de/e/cold-brew-anleitung}{Coffee Circle: Cold Brew Anleitung}},
    calory = {2kcal | 0g Protein | 0g KH | 0g Fett},
    price = {0,60€},
    created = {10.06.2026},
    %rating = {1},
]
{Cold Brew}

%\graph{% pictures
%    big=pic/getraenke/02_cold-brew
%}

\introduction{%
Cold Brew ist Kaffee für Leute, die Kaffee wollen, aber nicht direkt von Bitterstoffen angeschrien werden möchten. Fruchtig, kalt, sehr sommerlich und eigentlich lächerlich einfach.
}

\ingredients{
    100 g & Kaffee, grob gemahlen\index{Getränke!Kaffee}\\
    1 l & Kaltes Wasser\index{Getränke!Wasser}
}

\preparation{%
    \step%
    Kaffee grob mahlen und mit kaltem Wasser in einem Gefäß oder Cold-Brew-Maker vermischen.
    
    \step%
    Gut umrühren, damit das Kaffeepulver vollständig mit Wasser bedeckt ist.
    
    \step%
    Gefäß abdecken und mindestens 12 Stunden im Kühlschrank ziehen lassen.
    
    \step%
    Cold Brew durch ein feines Sieb, einen Papierfilter oder den Filtereinsatz des Cold-Brew-Makers filtrieren.
    
    \step%
    Pur auf Eis servieren oder nach Geschmack mit Milch, Pflanzendrink oder Saft mischen.
}

\suggestion[Konzentrat]{%
    Für ein Cold-Brew-Konzentrat 200 g Kaffee auf 1 l Wasser verwenden und später mit Wasser, Milch oder Pflanzendrink auf Trinkstärke verdünnen.
}

\hint{
    Grober Mahlgrad ist hier wichtig. Zu fein gemahlener Kaffee macht das Filtern unnötig nervig und kann den Cold Brew schnell bitterer machen.
}

\end{recipe}